\documentclass[10pt,twocolumn]{article}
\usepackage[T1]{fontenc}
\usepackage[utf8]{inputenc}
\usepackage{lmodern}
\usepackage[margin=1.8cm,top=2.4cm,bottom=2cm,columnsep=0.6cm]{geometry}
\usepackage{fancyhdr}
\usepackage{titlesec}
\usepackage{booktabs}
\usepackage{amsmath,amssymb,amsfonts}
\usepackage{microtype}
\usepackage{xcolor}
\usepackage{mdframed}
\usepackage{parskip}
\usepackage{enumitem}
\usepackage{hyperref}

\definecolor{rulered}{RGB}{180,30,30}
\definecolor{boxgray}{RGB}{245,245,245}
\definecolor{keywordgray}{RGB}{100,100,100}

\pagestyle{fancy}
\fancyhf{}
\fancyhead[L]{\textsc{\small Claude Mag}}
\fancyhead[C]{\small\textit{Machine Learning Research Digest}}
\fancyhead[R]{\small 2026-07-08}
\fancyfoot[C]{\small\thepage}
\renewcommand{\headrulewidth}{0.8pt}

\titleformat{\section}{\normalsize\bfseries\color{rulered}}{}{0pt}{}%
  [\vspace{-4pt}\color{rulered}\rule{\columnwidth}{0.4pt}]
\titlespacing{\section}{0pt}{8pt}{4pt}

\hypersetup{colorlinks=true,urlcolor=rulered,linkcolor=black}

\begin{document}
\setlength{\parindent}{0pt}
\setlength{\parskip}{4pt}

{\small\color{keywordgray}\textbf{Keywords:}
  score-based sampling \textbullet{}
  diffusion models \textbullet{}
  log-concave distributions \textbullet{}
  sampling complexity}\par\vspace{4pt}

{\LARGE\bfseries\raggedright
  High-Accuracy Sampling for Diffusion Models
  and Log-Concave Distributions}\par\vspace{4pt}

{\large\itshape\raggedright
  First-Order Rejection Sampling Achieves
  Polylogarithmic Step Complexity---An Exponential Improvement}\par\vspace{6pt}

{\small\textbf{By} Fan Chen, Sinho Chewi, Constantinos Daskalakis \&
  Alexander Rakhlin (MIT, Yale University)}\par\vspace{2pt}

{\small\color{keywordgray}arXiv:2602.01338 \textbullet{}
  ICML 2026 Outstanding Paper Award (Oral, Jul 8 2026)}
\par\vspace{2pt}\noindent{\color{rulered}\rule{\columnwidth}{0.6pt}}\vspace{6pt}

Every diffusion model sampler deployed in practice requires polynomially many
denoising steps to achieve $\delta$-accurate output. This paper ends that
era. The First-Order Rejection Sampling (FORS) algorithm attains $\delta$
accuracy in $\widetilde{O}(d_\star\,\mathrm{polylog}(1/\delta))$ steps ---
an \textbf{exponential improvement} over all prior work. The same construction
yields the first polylogarithmic sampler for general log-concave distributions
using only gradient evaluations.

\begin{mdframed}[backgroundcolor=boxgray,linecolor=rulered,linewidth=1pt,
    innerleftmargin=6pt,innerrightmargin=6pt,
    innertopmargin=6pt,innerbottommargin=6pt]
\textbf{\small AT A GLANCE}\par\vspace{4pt}
\begin{description}[leftmargin=0pt,labelsep=4pt,itemsep=2pt,parsep=0pt,
    font=\small\bfseries]
  \item[Problem:] All existing diffusion model samplers require
    $\mathrm{poly}(1/\delta)$ steps to reach $\delta$-accuracy; the gap to
    optimal has been an open problem.
  \item[Why it matters:] Production samplers operate at coarse $\delta$ to
    limit step count; fine-grained accuracy has been computationally
    prohibited. FORS removes this constraint.
  \item[Method:] First-Order Rejection Sampling (FORS), a meta-algorithm
    that simulates exact rejection sampling using only score (gradient)
    evaluations.
  \item[Result:] $\delta$-error in $\widetilde{O}(d_\star\,\mathrm{polylog}(1/\delta))$
    steps --- exponential improvement; also first polylog sampler for
    log-concave distributions.
  \item[Evidence:] Formal complexity bounds; comparison with Langevin,
    DDPM, exponential integrator baselines on synthetic log-concave targets.
\end{description}
\end{mdframed}

\section{The Big Picture}

Score-based generative models learn to reverse a diffusion process by
estimating $\nabla \log p_t(x)$ at each noise level $t$. At sampling time,
an algorithm uses these score estimates to walk from a Gaussian prior back
toward the data distribution. The quality of the resulting sample depends on
both how accurately scores are estimated and how finely the reverse-time SDE
is discretized.

For a decade, progress on the discretization side has been incremental:
better numerical integrators (Euler--Maruyama, Runge--Kutta, exponential
integrators) reduced the per-step discretization error but could not
escape the $\mathrm{poly}(1/\delta)$ barrier. This paper identifies the
barrier's root cause---the inability of discretization methods to correct
global integration error---and bypasses it entirely via rejection sampling.

\section{Why Existing Approaches Fall Short}

Discretization-based samplers approximate the reverse SDE as a sequence of
local updates. Each step commits to a trajectory segment; errors accumulate
across steps in a way that is fundamentally bounded by the product of
per-step error and step count. Halving the step size halves local error but
doubles step count, leading to the poly$(1/\delta)$ wall.

High-accuracy rejection sampling can escape this wall: it proposes from a
tractable distribution and accepts or rejects based on an exact density ratio,
which is globally correct regardless of the number of steps. The obstacle
has been that evaluating the exact density $p_t(x)$ is intractable for
learned diffusion models---one can evaluate $\nabla \log p_t$ (the score)
but not $p_t$ itself.

Prior high-accuracy methods (convex optimization over log-concave targets)
required density evaluations or were restricted to special distribution classes.
No general-purpose polylog sampler existed.

\section{Core Mechanism}

FORS shows that rejection sampling can be \emph{simulated} using only score
evaluations. The key insight is that the density ratio at the core of
rejection sampling can be expressed as a path integral of score differences,
which is estimable from accumulated score queries along a proposed trajectory.

Specifically, FORS proposes a candidate path from a Gaussian under score
guidance, then estimates the probability of accepting that candidate by
computing the Girsanov-weight of the path under the target reverse SDE. This
weight depends only on score evaluations at points along the path, not on
point-wise density values.

\section{Technical Formulation}

Let $p_0$ be the data distribution and $p_t = p_0 * \mathcal{N}(0, t^2 I)$.
Given score estimates $\tilde{s}_t$ with $L^2$ accuracy:
\[
\mathbb{E}\bigl[\|\tilde{s}_t(x) - \nabla\!\log p_t(x)\|^2\bigr]
  \leq \tilde{\delta}^2,
\]
the FORS acceptance weight for a proposed path $\{x_t\}_{t\in[0,T]}$ is
\[
w = \exp\!\left(
  \int_0^T \bigl(s_t(x_t) - \tilde{s}_t(x_t)\bigr)^{\!\top} dx_t
  - \tfrac{1}{2}\int_0^T \bigl(\|s_t\|^2 - \|\tilde{s}_t\|^2\bigr) dt
\right).
\]
The main theorem establishes that accepting paths proportional to $w$ yields
samples within total variation $\delta$ of $p_0$ in
\[
T = \widetilde{O}\!\left(d_\star \cdot \mathrm{polylog}(1/\delta)\right)
\]
score evaluations, where $d_\star$ is the intrinsic dimension of $p_0$.
Under a non-uniform $L$-Lipschitz score condition the bound tightens to
$\widetilde{O}(L\,\mathrm{polylog}(1/\delta))$.

\section{Learning or Inference Procedure}

\begin{enumerate}[itemsep=2pt,parsep=0pt]
  \item Initialize $x_T \sim \mathcal{N}(0, T^2 I)$ from the Gaussian prior.
  \item Simulate a reverse-time diffusion proposal path $\{x_t\}$ using
    score estimates $\tilde{s}_t$ (standard DDPM or exponential integrator step).
  \item Compute the FORS acceptance weight $w$ from score evaluations along
    the proposed path.
  \item Accept the terminal point $x_0$ with probability $\min(1, w/w_{\max})$.
  \item If rejected, restart from step 1. Expected restarts: $O(1)$.
  \item Return accepted $x_0$ as the sample.
\end{enumerate}

\section{What the Guarantee Says}

\textbf{Theorem (informal).} Given score estimates with $L^2$ error $\leq\tilde{\delta}$,
FORS returns a sample within total variation $\delta = O(\tilde{\delta})$
of $p_0$ using $\widetilde{O}(d_\star\,\mathrm{polylog}(1/\tilde{\delta}))$ score
evaluations. This is \emph{tight} up to polylog factors: any algorithm using
only $L^2$-accurate score estimates requires $\Omega(d_\star / \delta^2)$
steps (lower bound from information theory), but FORS achieves polylog ---
the gap arises because FORS uses the score path integrally, not pointwise.

\section{Experimental Findings}

On synthetic log-concave targets (Gaussian mixtures, strongly log-concave
distributions up to dimension 512):
\begin{itemize}[itemsep=2pt,parsep=0pt]
  \item FORS achieves TV$\leq 10^{-4}$ in $\approx 80$ score evaluations.
  \item DDPM requires $>50{,}000$ steps for the same accuracy.
  \item Exponential integrators (DPM-Solver++) require $>5{,}000$ steps.
  \item Step count grows as $O(\log^3(1/\delta))$ empirically, matching theory.
\end{itemize}
On image diffusion (score model for CIFAR-10), FORS attains FID comparable
to 1{,}000-step DDPM in 120 FORS steps (including restarts).

\section{Ablations and Interpretation}

\textbf{Score accuracy $\tilde{\delta}$:} Step count scales as
$O(\tilde{\delta}^{-2}\,\mathrm{polylog}(1/\delta))$, so coarser score
estimates increase cost modestly.

\textbf{Proposal distribution:} FORS tolerates diverse proposal distributions;
Euler--Maruyama proposals are sufficient. More accurate proposals reduce
expected restart count but do not affect the polylog exponent.

\textbf{Intrinsic dimension $d_\star$:} For image distributions, $d_\star \ll d$
(ambient dimension); FORS automatically exploits this, explaining why
the empirical step count on CIFAR-10 is far below the $d = 3072$ worst case.

\textbf{Practical implication:} The result motivates a new generation of
samplers that combine FORS's theoretical guarantee with learned proposal
distributions to minimize expected restart cost.

\section*{Reference}

Fan Chen, Sinho Chewi, Constantinos Daskalakis, Alexander Rakhlin.
\textbf{High-Accuracy Sampling for Diffusion Models and Log-Concave Distributions}.
arXiv:2602.01338, February 2026.
\url{https://arxiv.org/abs/2602.01338}

\end{document}
